% Chapter 10: Network Simulations — Meta-MAPG on Trade Networks

\chapter{Network Simulations: Meta-MAPG on Economic Trade Networks}
\label{ch:simulations}

The preceding chapters developed the Meta-MAPG theorem (Chapter~\ref{ch:meta-mapg}) as a theoretical framework for multi-agent learning under non-stationarity, and validated it through the lens of deployed autonomous agents (Section~\ref{sec:aoc-interpretation}). In this chapter, we instantiate the framework on a structured economic network---the West Africa Free Trade Zone---to test the theorem's predictions about cascade damping and emergent cooperation in a setting where the network topology explicitly mediates inter-agent coupling.

\section{The West Africa Trade Network as Stochastic Game}
\label{sec:trade-network-game}

We model the ECOWAS economic network as an $n$-agent stochastic game $\mathcal{M}_n$ (Definition~\ref{def:stochastic-game}), where each agent controls the trade policy of one country.

\paragraph{State space.} The state $s_t$ is the weighted multigraph $\mathcal{G}_t = (\mathcal{V}, \mathcal{E}_t, \mathbf{w}_t)$, where $\mathcal{V}$ comprises 45 West African cities, $\mathcal{E}_t$ is the set of typed edges (trade, financial, infrastructure, political, cultural, migratory, labour), and $\mathbf{w}_t$ contains edge-level attributes (trade volume, tariff rate, infrastructure quality). The state evolves dynamically as agents' actions modify edge weights and node attributes.

\paragraph{Agents and actions.} Agent~$i$ controls all cities belonging to country~$i$. The action space $\mathcal{A}^i$ includes:
\begin{itemize}
  \item \textsc{lower\_tariff} / \textsc{raise\_tariff}: modify tariff rates on trade edges
  \item \textsc{invest\_infrastructure}: reduce weight (improve connectivity) on infrastructure edges
  \item \textsc{form\_bilateral}: establish new bilateral trade agreements
  \item \textsc{exit\_bloc} / \textsc{rejoin\_bloc}: change ECOWAS membership status
  \item \textsc{no\_action}: maintain current policy
\end{itemize}

\paragraph{Transition dynamics.} The transition $\mathcal{P}(s_{t+1} \mid s_t, \mathbf{a}_t)$ incorporates two mechanisms:
\begin{enumerate}
  \item \emph{Direct effects}: each action modifies the relevant edge attributes (e.g., lowering tariffs increases trade volume on adjacent edges).
  \item \emph{Cascade effects}: disruptive actions (bloc exit) trigger cascading disruptions computed by the economic cascade simulator. When country~$i$ exits ECOWAS, all trade and financial edges to $i$'s cities are severed, potentially fragmenting the remaining network into disconnected components and isolating dependent cities.
\end{enumerate}

This cascade mechanism is the physical realisation of the inter-agent coupling that Term~3 of the Meta-MAPG gradient captures.

\paragraph{Reward.} Agent~$i$'s reward $R^i(s_t, \mathbf{a}_t)$ combines:
\begin{equation}
  R^i = \Delta\text{TradeVol}^i + \lambda_{\text{ftz}} \cdot \text{FTZScore}^i - \lambda_{\text{cascade}} \cdot \text{CascadePenalty}^i,
\end{equation}
where $\Delta\text{TradeVol}^i$ is the change in total trade volume for country~$i$'s cities, $\text{FTZScore}^i$ is the composite FTZ impact score (connectivity, port access, tariff exposure, diversification), and $\text{CascadePenalty}^i$ captures the indirect damage from network fragmentation.

\paragraph{Partial observability.} Each agent observes its own cities' features and direct neighbours' features, but not the full graph---consistent with real-world information asymmetries in trade policy. This maps to the POMDP framework of Chapter~\ref{ch:pgt-proof}.

\section{Network Topology and Agent Coupling}
\label{sec:network-coupling}

The key insight connecting the Meta-MAPG framework to the network infrastructure: \emph{the network topology determines the coupling structure between agents}.

\paragraph{Coupling strength.} For agents $i$ and $j$, define the coupling strength:
\begin{equation}\label{eq:coupling}
  \kappa_{ij} = \frac{\sum_{e \in \mathcal{E}^{ij}_{\text{trade}}} \text{vol}(e)}{\sum_{e \in \mathcal{E}_{\text{trade}}} \text{vol}(e)},
\end{equation}
where $\mathcal{E}^{ij}_{\text{trade}}$ is the set of trade edges between countries $i$ and $j$. This determines the magnitude of Term~3: if Nigeria and Ghana share \$2.3B in quarterly trade (high $\kappa$), Nigeria's tariff policy has a large effect on Ghana's learning environment.

\paragraph{Articulation points as critical agents.} Cities identified as articulation points in the graph have outsized influence: their country's exit would disconnect the ECOWAS subgraph. The agents controlling these cities have the largest Term~3 contributions, because their actions affect the most peers' environments.

\paragraph{Spectral radius and cascade threshold.} The spectral radius $\rho(\mathbf{A})$ of the adjacency matrix determines whether perturbations grow (cascade regime, $\rho > 1$) or decay. From Section~\ref{sec:aoc-interpretation}:
\begin{equation}
  \delta\bm{\phi}_\ell^{-i} \approx \prod_{k=0}^{\ell-1} \frac{\partial\bm{\phi}_{k+1}^{-i}}{\partial\bm{\phi}_k} \cdot \delta\bm{\phi}_0^i.
\end{equation}
When the spectral radius of these Jacobian products exceeds unity, perturbations propagate supercritically. An agent computing Term~3 modulates its policy to reduce this effective spectral radius.

\section{Experimental Design: Three-Way Ablation}
\label{sec:ablation}

We run three variants to isolate the contribution of each gradient term:

\begin{center}
\begin{tabular}{lcccl}
  \hline
  \textbf{Variant} & \textbf{Term 1} & \textbf{Term 2} & \textbf{Term 3} & \textbf{Equivalent to} \\
  \hline
  Independent PG & \checkmark & \texttimes & \texttimes & Standard REINFORCE \\
  Meta-PG        & \checkmark & \checkmark & \texttimes & \cite{alshedivat2018} \\
  Meta-MAPG      & \checkmark & \checkmark & \checkmark & Full Theorem~\ref{thm:meta-mapg} \\
  \hline
\end{tabular}
\end{center}

For each variant, we:
\begin{enumerate}
  \item Train for $M$ meta-steps with $L$ inner-loop updates per step, collecting $K$ trajectories per inner step.
  \item After training, evaluate over $E$ episodes measuring: average discounted return, cascade depth (number of cities disrupted per exit event), cascade severity (fraction of total ECOWAS trade affected), and exit frequency (how often agents choose \textsc{exit\_bloc}).
\end{enumerate}

\section{Cascade Damping Prediction}
\label{sec:cascade-damping}

The central testable prediction: under Meta-MAPG, the expected cascade depth is bounded by

\begin{equation}\label{eq:cascade-bound}
  \E\!\left[\text{cascade\_depth} \mid \text{Meta-MAPG}\right] \leq \frac{\E\!\left[\text{cascade\_depth} \mid \text{independent}\right]}{1 + \alpha \cdot L \cdot |\mathcal{E}|},
\end{equation}

where $\alpha$ is the inner-loop learning rate, $L$ is the number of inner steps, and $|\mathcal{E}|$ is the number of edges in the trade network.

\paragraph{Mechanism.} Term~3 creates $O(L \cdot |\mathcal{E}|)$ additional gradient signal that penalises actions causing large perturbations to peers' environments. On the trade network, this means agents learn to avoid bloc exits (which trigger cascades) and prefer cooperative actions (tariff reduction, infrastructure investment) that improve the shared environment.

\paragraph{Network topology modulates the bound.} The damping is strongest on dense networks with high coupling $\kappa_{ij}$, and weakest on sparse networks where agents are nearly independent. The West Africa graph, with its hub-and-spoke structure centred on Nigeria and C\^ote d'Ivoire, provides a realistic intermediate case.

\section{Connection to the GNN-TCN Prediction Model}
\label{sec:gnn-connection}

The GNN-TCN model (spatial graph attention + temporal convolution) trained on the same West Africa network provides a complementary perspective:

\begin{itemize}
  \item \textbf{GAT layer} = spatial encoding of agent coupling. The attention weights learned by the Graph Attention Network correspond to the coupling strengths $\kappa_{ij}$ that determine Term~3's magnitude. High attention weight between Nigeria and Ghana means the GNN has learned that these economies are tightly coupled---the same structural fact that makes Term~3 large.
  \item \textbf{TCN layer} = temporal encoding of policy chains. The dilated causal convolutions capture multi-step dependencies across quarters---the temporal counterpart of the policy chain $\bm{\phi}_0 \to \bm{\phi}_1 \to \cdots \to \bm{\phi}_L$ in the inner loop.
  \item \textbf{Prediction targets} = FTZ impact scores for target cities. These are exactly the quantities that Meta-MAPG agents should learn to optimise: trade volume growth for their own cities, mediated by the network structure.
\end{itemize}

The GNN-TCN model can serve as a \emph{world model} for the philosopher-king layer: given a proposed joint action (set of tariff changes, infrastructure investments, bloc membership changes), the model predicts the economic impact on each city, enabling counterfactual reasoning (Pearl's Level~3) without running the full stochastic game.

\section{From Agents of Chaos to Trade Networks}
\label{sec:aoc-to-network}

The Agents of Chaos case studies (Section~\ref{sec:aoc-interpretation}) and the trade network simulations test the same theoretical prediction---Term~3 damping---in two complementary settings:

\begin{center}
\begin{tabular}{lll}
  \hline
  & \textbf{Agents of Chaos} & \textbf{Trade Network} \\
  \hline
  Environment & Discord server + tools & ECOWAS multigraph \\
  Agents & 6 LLM agents & $n$ country-level policy agents \\
  Coupling mechanism & Shared communication channel & Shared trade edges \\
  Cascade trigger & Spoofed identity broadcast & Bloc exit \\
  Term 3 success & CS16: emergent safety coordination & Reduced exit frequency \\
  Term 3 failure & CS11: uncontrolled cascade & Network fragmentation \\
  Observable & Message propagation depth & Cascade depth + severity \\
  \hline
\end{tabular}
\end{center}

Both settings confirm the core prediction: agents that model peer learning (Term~3) produce damped cascades compared to independent learners. The trade network additionally provides \emph{quantitative} predictions via the cascade bound~\eqref{eq:cascade-bound} and the spectral analysis of the adjacency matrix.

\section{Implementation}
\label{sec:implementation}

The simulation code is implemented in \texttt{west\_africa/bridge/}, reconciling the MARL framework with the network infrastructure:

\begin{itemize}
  \item \texttt{marl\_network\_env.py}: wraps \texttt{WestAfricaGraph} as a stochastic game (Definition~\ref{def:stochastic-game}), implementing the action effects and cascade transitions.
  \item \texttt{meta\_mapg.py}: implements the three-term gradient computation from Theorem~\ref{thm:meta-mapg}, with ablation support for isolating Terms~1, 2, and 3. The peer-learning Jacobian (Term~3) uses the network coupling strength $\kappa_{ij}$ from~\eqref{eq:coupling}.
  \item \texttt{cascade\_damping.py}: runs the ablation experiments and tests the cascade bound~\eqref{eq:cascade-bound}, connecting the cascade simulator to the MARL training loop.
\end{itemize}

The \texttt{bridge} module unifies code that was previously separate: the economic network analysis (graph metrics, cascade simulation, GNN-TCN predictions) and the MARL theory (stochastic games, policy gradients, meta-learning). The network topology enters the gradient computation through $\kappa_{ij}$, and the gradient computation produces policies that reshape the network through their economic actions.
